\section{导论：Beyond LLMs 的两层含义}

这堂 instructors lecture 把两个看似分散的问题放在一起。第一层是模型内部：能力为什么会随规模、提示和中间表示而改变，怎样用 Chain-of-Thought、search、program execution 与 data efficiency 获得更强推理。第二层是模型外部：一个语言模型怎样成为 agent system 的计算核心，并与 memory、tools、environment feedback、permissions 和 other agents 共同完成长期任务。

\lecturefigure{slide-01-title.jpg}{CS25 V3 instructors lecture 的四个主题}{Stanford 课堂视频 00:00:21}

\begin{knowledgebox}{读图：四个主题构成一条能力扩展链}
Emergent abilities 讨论规模与测量；Intermediate-Guided Reasoning 讨论怎样组织推理过程；BabyLM 反问是否必须依靠海量数据；Agents 则把模型放进能观察、行动、记忆和纠错的系统。它们共同回答“怎样超越一次前向生成”，而不是四段互不相关的热点综述。
\end{knowledgebox}

\begin{importantbox}{本讲的统一系统视角}
可把能力写成三层组合：
\[
\text{Observed capability}
=\text{Model}
\circ\text{Inference procedure}
\circ\text{Environment loop}.
\]
Model 决定参数化知识与表征，inference procedure 决定 prompt、search 或 program execution，environment loop 决定工具、反馈、记忆和权限。课堂后半段的 agent 正是把第三层显式化。
\end{importantbox}

\subsection{证据边界}

前半段主要是论文导读，后半段包含 MultiOn 原型演示和架构设想。讲义会把论文结论、课堂解释、原型 capability claim 与工程推论分开；尤其不能把 2023 年的飞行预订、驾驶测试或 universal action API 演示写成当前产品规格或通用可靠性证明。

\subsection{本章小结}

Beyond LLMs 不是离开语言模型，而是同时改变模型规模、推理过程和系统边界。接下来先从最容易被误解的 emergence 开始，再逐步走向结构化推理、样本效率和 agent architecture。

